\documentclass[book.tex]{subfiles}

\begin{document}

\chapter{Molecular Dynamics}
\label{chap:molecular_dynamics}
\noindent\rule{11cm}{0.4pt} \\
\textbf{Prerequisites:} \autoref{chap:newtons_laws}, \autoref{chap:electrostatics}  \\
\textbf{Difficulty Level:} ** \\
\noindent\rule{11cm}{0.4pt}
\vspace{5mm}

Molecular dynamics methods numerically solve Newton's equations for a system of interacting particles. Molecular dynamics methods offer the powerful advantage of efficiency compared with quantum mechanical simulation methods, but at the cost of modeling accuracy.

\section{The Basics of Molecular Dynamics}

At its core, molecular dynamics is roughly a variant of Euler's method for solving systems of ODEs. The gradient of the energy (the force) is used to update the positions of all atoms in the system repeatedly. This basic formula is modified to enforce boundary conditions, and temperature and pressure controls.

\begin{figure}
    \centering
    % \includegraphics{figures/part2b/molecular_dynamics/molecular_dynamics_algorithm.png}
    \includegraphics{figures/part2b/molecular_dynamics/MD_algorithm.png}
    \caption{Conceptual pseudocode of molecular dynamics simulation algorithm}% adapted from \url{https://en.wikipedia.org/wiki/Molecular_dynamics}. TODO: Make a new figure of our own}
    \label{fig:my_label}
\end{figure}

We can define the basic molecular dynamics algorithm in Physika as follows

\begin{lstlisting}[language=python]
class MolecularDynamics($r: \mathbb{R}^{3N}$, $v: \mathbb{R}^{3N}$, $\Delta t: \mathbb{R}$):

  def $\lambda$(N : $\mathbb{N}$): $(\mathbb{R}^{3N}, \mathbb{R}^{3N})$:
    for _ in range(N):
      $r$ = $r$ + $v \Delta t$ + $\frac{1}{2} \alpha \Delta t^2$
      $F$ = $-\nabla$ ForceField($r$)
      $r$ = $r$ + ForceUpdate(F, $\Delta t$)
      $v$ = $v$ + VelocityUpdate(F, $\Delta t$)
    return (r, v)
\end{lstlisting}

\section{Cleanup}

One of the most important practical steps in molecular dynamics is cleaning up input protein structure files. Common cleanup steps include fixing protonation, creating disulfide bonds, and placing side-chains correctly. 

\section{Force Fields}

Force fields are the technical term for the classical potential functions that are used in molecular dynamics systems to model interesting physics. Force fields typically correspond to classical energies are are made up of constituent energy terms. A common such term is the Lennard-Jones potential

\begin{align}
    U(r) &= 4 \epsilon \left [ \left ( \frac{\sigma}{r} \right )^{12} + \left( \frac{\sigma}{r} \right )^6 \right ]
\end{align}

The Lennard-Jones potential provides an approximate model for Van der Waals interactions 

%\textcolor{red}{TODO: Add a derivation of Van der Waals}

%\textcolor{red}{TODO: We describe Van der Waals somewhere in the stat-mech section. Cross link to here.}

Force fields are typically parameterized by a collection of tunable parameters. Traditionally, force field parameters were manually tuned but as we will learn in later chapters, force fields are increasingly learned using the techniques of machine learning and differentiable programming.

\section{Neighbor Lists}

One of the most important steps in a molecular dynamics update operation is computing the neighbors of a given atom which contribute to force estimates. The set of current neighbors of an atom is maintained by a data structure known as a neighbor list. Any molecular dynamics code must implement a neighbor list.

\section{Ewald Summation}

The most inefficient step in the computation of a force field is computing pairwise interaction terms which requires $O(n^2)$ computations. Ewald summation provides a method for reducing the computational cost of the force computation to $O(n \log n)$.

\section{Stability of Integration}


One of the challenges of molecular dynamics, especially with learned force fields, is stability of integration. Integration can cause a system to become unstable and "blow apart" in simulation if the force field isn't well calibrated. Heuristically designed force fields are known to be well calibrated in practice, but learned force fields are still unstable. There are some signs that this instability may be due to the numerical instability of using gradient information across very long timescales.

%\textcolor{red}{TODO: Add a mathematical description of Ewald summation.}

\end{document}

